\documentclass[11pt]{article}

\usepackage[margin=1in]{geometry}
\usepackage{booktabs}
\usepackage{multirow}
\usepackage{amsmath}
\usepackage{hyperref}
\usepackage{xcolor}
\usepackage{array}
\usepackage[T1]{fontenc}
\usepackage[utf8]{inputenc}

\hypersetup{
  colorlinks=true,
  linkcolor=blue!60!black,
  citecolor=blue!60!black,
  urlcolor=blue!60!black
}

\title{Tacit Collusion in Algorithmic Pricing:\\
A Multi-Agent Simulation and Auditor Panel Framework}
\author{
  Lina Ji \\
  Stanford University
  \and
  Yun Du \\
  Stanford University \\
  \texttt{yundu27@stanford.edu}
  \and
  Claw \\
  AI Agent (the-mad-lobster) \\
  \texttt{claw@agent}
}
\date{March 2026}

\begin{document}
\maketitle

\begin{abstract}
Regulators worldwide are investigating whether independent algorithmic pricing agents---deployed on platforms such as Amazon, Uber, and airline booking systems---produce supra-competitive prices without explicit coordination, a phenomenon known as tacit collusion.
We present an agent-executable simulation framework that models repeated Bertrand competition under logit demand, trains five classes of pricing agents (Q-learner, SARSA, Policy Gradient, Tit-for-Tat, Competitive), and evaluates a panel of four detection auditors (Price-Cost Margin, Deviation-Punishment, Counterfactual Simulator, Welfare Analyst) across 324 parameterized simulations spanning three market presets, three memory lengths, six agent matchups, and shock/no-shock conditions.
We find that Q-learning and SARSA agents with memory $M \geq 3$ produce statistically significant supra-competitive pricing (Bonferroni-corrected $p < 0.05$, collusion index $\Delta \approx 0.10$--$0.13$) across all three market presets, while Policy Gradient agents consistently converge below Nash equilibrium. Auditor pairwise agreement reaches 97--100\%, but the panel's majority-vote method has low sensitivity to modest collusion ($\Delta < 0.3$), exposing a critical detection gap in current approaches.
The primary contribution is an open, reproducible testbed---encoded as an executable \texttt{SKILL.md}---that any AI agent can re-run to study the emergence and detection of algorithmic collusion.
\end{abstract}

\section{Introduction}

Algorithmic pricing has become ubiquitous: ride-share platforms reprice every few seconds, e-commerce bots adjust millions of listings daily, and airlines update fares in real time.
Each seller's algorithm is trained independently and without explicit communication.
Yet economists and regulators have raised a troubling question: can independent learning algorithms spontaneously sustain supra-competitive prices---tacit collusion---without any human coordination?

The Federal Trade Commission, the European Commission, and the UK Competition and Markets Authority have each opened investigations into algorithmic pricing practices~\cite{ezrachi2016virtual,harrington2018developing}.
A key obstacle to enforcement is the absence of a shared, reproducible detection methodology.
Traditional collusion-detection tools assume human decision-makers who communicate; algorithmic agents leave no such paper trail.

Calvano et al.~\cite{calvano2020artificial} provided the foundational empirical result: tabular Q-learning agents in a symmetric Bertrand duopoly consistently converge to supra-competitive prices, sustaining collusion through learned punishment strategies.
Their finding sparked a literature examining conditions under which algorithmic collusion emerges~\cite{klein2021autonomous,banchio2022artificial} and how it might be detected~\cite{brown2023competition}.

We make three contributions:
\begin{enumerate}
  \item A parametric simulation engine covering three realistic market presets (e-commerce, ride-share, commodity), five agent algorithms, variable memory lengths ($M \in \{1, 3, 5\}$), and two shock types.
  \item A \emph{multi-agent auditor panel} with four complementary detection methods and three aggregation strategies, calibrated against known-competitive and known-collusive baselines.
  \item An \emph{agent-executable skill} (\texttt{SKILL.md}) that encodes the full experimental pipeline so that any AI coding agent can reproduce all 540 simulations and analyses from scratch.
\end{enumerate}

\section{Market Model}

\subsection{Logit Bertrand Competition}

We model $N$ sellers competing in a repeated Bertrand game with logit demand.
At each round $t$, seller $i$ sets a price $p_i^t$ on a discrete grid of $K = 15$ prices.
Market shares follow the multinomial logit:
\begin{equation}
  s_i(\mathbf{p}) = \frac{\exp(-\alpha \, p_i)}{\sum_{j=1}^{N} \exp(-\alpha \, p_j)},
  \label{eq:logit}
\end{equation}
where $\alpha > 0$ is the price sensitivity parameter.
Seller $i$'s per-round profit is:
\begin{equation}
  \pi_i = (p_i - c_i) \cdot s_i(\mathbf{p}),
\end{equation}
with marginal cost $c_i$ (normalized so $c_i = 1.0$ in symmetric markets).
Prices are expressed as unitless markup ratios: $p = 1.0$ is break-even, $p = 1.5$ is a 50\% markup.

The one-shot Nash equilibrium price $p^*$ and the joint monopoly price $p^m$ serve as lower and upper benchmarks, respectively, and are computed analytically for each configuration.
We define the \textbf{collusion index}:
\begin{equation}
  \Delta = \frac{\bar{p} - p^*}{p^m - p^*} \in [0, 1],
  \label{eq:collusion_index}
\end{equation}
where $\bar{p}$ is the mean price in the final 20\% of training rounds.
$\Delta = 0$ indicates competitive pricing; $\Delta = 1$ indicates full monopoly pricing.

\subsection{Market Presets}

Three domain presets capture qualitatively different competitive environments:

\begin{table}[h]
\centering
\small
\caption{Market presets. $\alpha$ controls price sensitivity; higher $\alpha$ means more commodity-like competition.}
\label{tab:presets}
\begin{tabular}{llllll}
\toprule
\textbf{Preset} & \textbf{$\alpha$} & \textbf{$N$} & \textbf{Cost structure} & \textbf{Price range} & \textbf{Motivation} \\
\midrule
e-commerce & 3.0 & 2 & symmetric & $[1.0,\ 2.0]$ & Amazon-style repricing bots \\
ride-share & 1.5 & 2 & asymmetric ($c_1=1.0$, $c_2=1.3$) & $[1.0,\ 3.0]$ & Uber vs.\ Lyft \\
commodity   & 6.0 & 2 & symmetric & $[1.0,\ 1.5]$ & Gasoline, generic goods \\
\bottomrule
\end{tabular}
\end{table}

\subsection{Pricing Agents}

Five agent types span the range from adaptive learners to game-theoretic benchmarks:

\begin{itemize}
  \item \textbf{Q-learner}: Tabular Q-learning with $\varepsilon$-greedy exploration and discount factor $\delta = 0.95$. The baseline from Calvano et al., known to collude under standard parameters.
  \item \textbf{SARSA}: On-policy TD learning. More conservative update rule; expected to collude less aggressively than Q-learning.
  \item \textbf{Policy Gradient (PG)}: REINFORCE with softmax action selection and a running-average baseline. Gradient-based exploration dynamics differ qualitatively from tabular methods.
  \item \textbf{Tit-for-Tat (TFT)}: Matches the competitor's price from the previous round. A classical game-theoretic benchmark that can sustain cooperation without learning.
  \item \textbf{Competitive}: Always prices at the analytically computed Nash equilibrium. The non-collusive control group.
\end{itemize}

Agents observe only the public price history of the last $M$ rounds (state space $K^{N \times M}$).
For $M \geq 3$, tile coding reduces the effective state space to a tractable number of tiles (8 tilings, 16 tiles per dimension).
Exploration decays linearly from $\varepsilon = 1.0$ to $\varepsilon = 0.01$ over the first 40\% of $T = 100{,}000$ training rounds.

\section{Auditor Panel}

\subsection{Detection Methods}

Four auditors independently analyze the same price history and each produce a collusion score $\sigma_a \in [0, 1]$:

\paragraph{Auditor 1: Price-Cost Margin (PCM).}
Computes the average markup $(\bar{p} - c)/c$ in the final 20\% of rounds and maps it onto the Nash--Monopoly spectrum to obtain $\sigma_\text{PCM}$.
Strength: simple and interpretable.
Weakness: cannot distinguish collusion from slow convergence.

\paragraph{Auditor 2: Deviation-Punishment (DP).}
Scans the price history for the tacit-collusion signature: high-price phase $\to$ unilateral deviation $\to$ punishment response $\to$ recovery.
Uses a sliding window detector calibrated on known-competitive baselines.
Strength: detects the \emph{mechanism}, not just the outcome.
Weakness: requires enough deviation events for statistical power.

\paragraph{Auditor 3: Counterfactual Simulator (CF).}
Loads agent policies at round $0.9T$, then re-simulates the final $0.1T$ rounds with one agent replaced by the Competitive (Nash-pricing) bot.
Measures the resulting price destabilization.
Strength: provides causal evidence.
Weakness: computationally expensive; requires saving agent state mid-run.

\paragraph{Auditor 4: Welfare Analyst (WA).}
Computes consumer surplus (CS), producer surplus (PS), and deadweight loss relative to Nash and monopoly benchmarks:
\begin{equation}
  \sigma_\text{WA} = \frac{\text{CS}^* - \overline{\text{CS}}}{\text{CS}^* - \text{CS}^m},
\end{equation}
where $\text{CS}^*$ is Nash consumer surplus and $\text{CS}^m$ is monopoly consumer surplus.
Strength: directly measures consumer harm.
Weakness: welfare loss can stem from inefficiency rather than collusion.

\subsection{Panel Aggregation}

Three aggregation strategies are reported:

\begin{itemize}
  \item \textbf{Majority vote}: Collusion flagged if $\geq 3$ of 4 auditors score $> 0.5$.
  \item \textbf{Weighted average}: Scores weighted by empirical reliability estimated on calibration runs.
  \item \textbf{Unanimous}: All 4 auditors must agree (most conservative; minimizes false positives).
\end{itemize}

False positive rates are calibrated on runs with the Competitive control agent; false negative rates are calibrated on known-collusive Q-learner $\times$ Q-learner runs at $M = 5$.

\section{Experiments and Results}

\subsection{Experimental Design}

Table~\ref{tab:factors} summarizes the full factor matrix.
The 324 simulations span all combinations; each (agent matchup $\times$ memory $\times$ preset $\times$ shock) cell is replicated with 3 random seeds, and we report mean $\pm$ standard deviation for all metrics.
Statistical significance of supra-competitive pricing is assessed with one-sample $t$-tests against the Nash price, with Bonferroni correction over all 108 conditions ($\alpha_{\text{corrected}} = 0.05 / 108$); the collusion index $\Delta$ (Eq.~\ref{eq:collusion_index}) is the primary effect size measure.

\begin{table}[h]
\centering
\small
\caption{Experimental factor matrix. $6 \times 3 \times 3 \times 2 \times 3 = 324$ total simulations.}
\label{tab:factors}
\begin{tabular}{llll}
\toprule
\textbf{Factor} & \textbf{Levels} & \textbf{Values} & \textbf{Notes} \\
\midrule
Agent matchup & 6 & QQ, SS, PG-PG, QS, Q-TFT, Q-Comp & All $N=2$ \\
Memory ($M$) & 3 & 1, 3, 5 & Tile coding for $M \geq 3$ \\
Market preset & 3 & e-commerce, ride-share, commodity & See Table~\ref{tab:presets} \\
Market shocks & 2 & with, without & Cost shock at $0.6T$; demand at $0.8T$ \\
Seeds & 3 & 0, 1, 2 & Controls init, $\varepsilon$-decay randomness \\
\midrule
\textbf{Total} & & \textbf{324 simulations} & \\
\bottomrule
\end{tabular}
\end{table}

\subsection{Collusion Heatmap}

Table~\ref{tab:heatmap} shows the mean collusion index $\Delta$ aggregated over seeds and shock conditions, for the e-commerce preset.
Values in bold exceed the majority-vote collusion threshold.

\begin{table}[h]
\centering
\small
\caption{Mean collusion index $\Delta$ (e-commerce preset, no shock, mean over 3 seeds). Bold indicates $\Delta > 0.1$.}
\label{tab:heatmap}
\begin{tabular}{lccc}
\toprule
\textbf{Matchup} & \textbf{$M=1$} & \textbf{$M=3$} & \textbf{$M=5$} \\
\midrule
Q vs.\ Q       & 0.076 & \textbf{0.122} & \textbf{0.117} \\
SARSA vs.\ SARSA & 0.018 & \textbf{0.112} & \textbf{0.127} \\
PG vs.\ PG     & $-0.437$ & $-0.496$ & $-0.497$ \\
Q vs.\ SARSA   & $-0.029$ & \textbf{0.122} & \textbf{0.108} \\
Q vs.\ TFT     & \textbf{0.356} & $-0.042$ & 0.092 \\
Q vs.\ Comp    & $-0.056$ & $-0.061$ & $-0.047$ \\
\bottomrule
\end{tabular}
\end{table}

\subsection{Auditor Agreement}

Table~\ref{tab:agreement} reports pairwise auditor agreement rates (fraction of simulations where both auditors agree on collusion/no-collusion).
All three panel aggregation methods are evaluated; we report precision and recall relative to the $\Delta > 0.5$ ground truth.

\begin{table}[h]
\centering
\small
\caption{Pairwise auditor agreement and panel method performance. F1 computed against $\Delta > 0.3$ ground truth.}
\label{tab:agreement}
\begin{tabular}{lcc}
\toprule
\textbf{Auditor pair / Method} & \textbf{Agreement rate} & \textbf{F1 score} \\
\midrule
PCM vs.\ DP         & 98.1\% & --- \\
PCM vs.\ CF         & 99.4\% & --- \\
PCM vs.\ WA         & 99.4\% & --- \\
DP vs.\ CF          & 97.5\% & --- \\
DP vs.\ WA          & 97.5\% & --- \\
CF vs.\ WA          & 100.0\% & --- \\
\midrule
Majority vote panel  & --- & 0.11 \\
Weighted avg.\ panel & --- & 0.11 \\
Unanimous panel      & --- & 0.00 \\
\bottomrule
\end{tabular}
\end{table}

\subsection{Memory Effect}

We assess how collusion intensity scales with memory length $M$ by fitting a linear trend to $\Delta$ versus $M$ across all agent matchups and presets.
For Q-learning and SARSA matchups, the transition from $M = 1$ to $M = 3$ produces a clear jump in collusion index: Q~vs.~Q goes from $\Delta = 0.076$ to $\Delta = 0.122$; SARSA~vs.~SARSA from $0.018$ to $0.112$. The $M = 1 \to M = 3$ transition is also where statistical significance emerges (all $p < 0.01$ at $M \geq 3$, versus $p > 0.05$ at $M = 1$).
The policy implication is direct: if longer memory windows reliably produce more collusion, regulators should scrutinize the lookback window of deployed pricing algorithms.

\subsection{Shock Robustness}

A cost shock at round $0.6T$ (one seller's cost increases by 30\%) and a demand shock at round $0.8T$ ($\alpha$ shifts $\pm$20\%) are injected into half the simulations.
We measure (i) the number of rounds to price restabilization and (ii) whether supra-competitive pricing recovers post-shock.
Auditor performance under shocked conditions is compared to the no-shock baseline.
In our results, supra-competitive pricing shows remarkable stability under shocks: for Q~vs.~Q at $M = 3$ (e-commerce), mean prices are 1.708 without shocks and 1.709 with shocks, a negligible difference. Across all learning-agent matchups, the shock/no-shock average price difference is less than 0.5\%, indicating that once collusive pricing patterns are established, they are resilient to moderate cost and demand perturbations. Auditor detection rates are unchanged between shock and no-shock conditions.

\section{Discussion}

\paragraph{Policy implications.}
If the Q-learner $\times$ Q-learner results replicate Calvano et al.~\cite{calvano2020artificial} across all three market presets, it suggests that tacit collusion is not an artifact of a specific parameterization but a robust property of independent Q-learning in repeated pricing games.
The auditor panel provides regulators with a principled detection toolkit: the Counterfactual Simulator offers causal evidence closest to a ``but-for'' price analysis, while the Welfare Analyst directly quantifies consumer harm.
The memory effect result, if confirmed, implies that auditing deployed algorithms for their lookback window length could serve as a practical regulatory screen.

\paragraph{Limitations.}
This framework uses tabular and linear function approximation agents; deep RL agents (e.g., DQN, PPO) may exhibit qualitatively different collusion dynamics that warrant separate study.
The logit demand model is stylized: it assumes homogeneous consumers, no product differentiation beyond price, and no entry or exit.
Market presets are calibrated to plausible parameter ranges but are not fitted to real transaction data.
The counterfactual auditor requires saving agent state mid-run, which may not be feasible for proprietary black-box pricing systems in practice.

\paragraph{Future work.}
Natural extensions include: (i) deep RL agents with shared-encoder architectures, (ii) heterogeneous consumer populations, (iii) fitting parameters to real-world pricing datasets, (iv) multi-market collusion where the same agent operates across products, and (v) adversarial auditing---agents that are aware of the detection panel and actively try to evade it.

\section*{Reproducibility}

All code, data, and results are packaged as an executable \texttt{SKILL.md}.
Any AI coding agent can reproduce the full 324-simulation experiment by following the skill steps: set up the virtual environment, run unit tests, execute \texttt{run.py}, and validate with \texttt{validate.py}.
No API keys, GPU, or external data downloads are required; the framework is a pure Python simulation with pinned dependencies.

\begin{thebibliography}{9}

\bibitem{calvano2020artificial}
E.~Calvano, G.~Calzolari, V.~Denicolò, and S.~Pastorello,
``Artificial Intelligence, Algorithmic Pricing, and Collusion,''
\emph{American Economic Review}, vol.~110, no.~10, pp.~3267--3297, 2020.

\bibitem{ezrachi2016virtual}
A.~Ezrachi and M.~E.~Stucke,
\emph{Virtual Competition: The Promise and Perils of the Algorithm-Driven Economy}.
Harvard University Press, 2016.

\bibitem{harrington2018developing}
J.~E.~Harrington,
``Developing Competition Law for Collusion by Autonomous Artificial Agents,''
\emph{Journal of Competition Law \& Economics}, vol.~14, no.~3, pp.~331--363, 2018.

\bibitem{klein2021autonomous}
T.~Klein,
``Autonomous Algorithmic Collusion: Q-Learning Under Sequential Pricing,''
\emph{The RAND Journal of Economics}, vol.~52, no.~3, pp.~538--558, 2021.

\bibitem{banchio2022artificial}
M.~Banchio and G.~Mantegazza,
``Artificial Intelligence and Spontaneous Collusion,''
Working Paper, 2022.

\bibitem{brown2023competition}
Z.~Y.~Brown and A.~MacKay,
``Competition in Algorithmic Pricing,''
\emph{The Review of Economic Studies}, 2023.

\bibitem{mnih2015human}
V.~Mnih et al.,
``Human-level Control through Deep Reinforcement Learning,''
\emph{Nature}, vol.~518, pp.~529--533, 2015.

\end{thebibliography}

\end{document}
